\subsection{不同避孕方式对女性性欲与性体验的影响}

避孕方式的选择不仅关乎"会不会怀孕"，也会或多或少影响\textbf{性欲、润滑、性交疼痛与性满意度}。没有一种方式对所有人都最好，关键是了解差异，选择与自己身体和生活方式最匹配的那一种。

\begin{itemize}
	\item \textbf{复方口服避孕药}：影响因人而异。部分女性报告性欲下降，可能与\textbf{性激素结合球蛋白（SHBG）升高、游离睾酮下降}有关；也有女性因不再担心怀孕而性欲提升。若怀疑与之相关，可与医生讨论更换剂型或改用非激素方式（详见上文"口服避孕药与性欲"）。
	\item \textbf{单纯孕激素避孕（皮下埋植、避孕针、孕激素宫内系统）}：部分人出现不规则出血、情绪波动或性欲变化；含左炔诺孕酮的宫内系统还常使月经量减少，对痛经与经期性交不适者可能是\textbf{加分项}。
	\item \textbf{含铜宫内节育器}：不含激素，一般不影响性欲；少数人出现经量增多、经期延长，可能间接影响经期前后的性体验。
	\item \textbf{屏障法（避孕套、避孕膜、宫颈帽）}：不影响激素；但避孕膜、宫颈帽需放置，部分人觉得打断节奏；乳胶过敏者可改用聚氨酯材质。
	\item \textbf{绝育手术（输卵管结扎）}：不影响激素分泌，理论上不影响性欲；部分女性因"彻底不用担心怀孕"而性满意度提升。
	\item \textbf{自然避孕法（安全期、体温法）}：不干扰身体，但需要在易孕期为避孕而禁欲或使用屏障法，对性生活的自发性与心理压力影响较大。
\end{itemize}

\noindent\textbf{实用建议}

\begin{itemize}
	\item \textbf{记录观察}：更换避孕方式后，若性欲、情绪、润滑或疼痛出现明显变化，记 2--3 个月再与医生讨论，避免只凭一两次感受下结论。
	\item \textbf{别自行停药}：无论怀疑哪种方式影响了性体验，都应\textbf{先与医生沟通再调整}，以免意外怀孕。
	\item \textbf{疼痛与干涩要单独处理}：绝经前后或哺乳期的干涩、性交痛，可借助水基/硅基润滑剂、局部雌激素（需医生处方）等方法改善，不必因此放弃满意的性生活。
\end{itemize}

\begin{tcolorbox}[colback=pink!5!white,colframe=red!50!black,title=贴心小叮咛]
避孕方式与性体验的关系是\textbf{双向}的：合适的避孕让人更放松、更享受，而不合适的方式可能悄悄压低性欲。若你在更换避孕方式后感到"不像自己了"，这不是矫情——把它作为一个真实的医学问题，带去和医生谈。
\end{tcolorbox}
